% ╔══════════════════════════════════════════════════════════════════════╗
% ║  LICEO V.A.L.  –  MÓDULO DE AUTOAPRENDIZAJE                        ║
% ║  Código: 06M-17 FL  |  Grado: 6°  |  Ruta: Formación Libre         ║
% ║  Unidad 17: Estadística. Lectura de tablas, gráficas y pictogramas ║
% ╚══════════════════════════════════════════════════════════════════════╝
\documentclass[12pt,letterpaper]{article}

% ── Codificación y tipografía ──────────────────────────────────────────
\usepackage{lmodern}
\usepackage[utf8]{inputenc}
\usepackage[T1]{fontenc}
\usepackage[spanish, es-noshorthands]{babel}
\usepackage{helvet}
\renewcommand{\familydefault}{\sfdefault}

% ── Geometría ─────────────────────────────────────────────────────────
\usepackage[letterpaper,left=1.5cm,right=1.5cm,top=1.8cm,bottom=1.8cm]{geometry}

% ── Gráficos y color ──────────────────────────────────────────────────
\usepackage[table]{xcolor}
\usepackage{graphicx}
\usepackage{eso-pic}
\usepackage{pifont}

% ── TikZ ──────────────────────────────────────────────────────────────
\usepackage{tikz}
\usetikzlibrary{babel,shapes.geometric,arrows.meta,calc,positioning,backgrounds,fit}

% ── Tcolorbox ─────────────────────────────────────────────────────────
\usepackage{tcolorbox}
\tcbuselibrary{most,skins,breakable}

% ── Paquetes adicionales ──────────────────────────────────────────────
\usepackage{array,tabularx,booktabs,longtable}
\usepackage{amsmath,amssymb}
\usepackage{enumitem}
\usepackage{multicol}
\usepackage{needspace}

% ═══════════════════════════════════════════════════════════════════════
%  PALETA INSTITUCIONAL VAL
% ═══════════════════════════════════════════════════════════════════════
\definecolor{valblue}       {RGB}{  0, 84,166}
\definecolor{valgreen}      {RGB}{  0,150, 64}
\definecolor{valred}        {RGB}{200, 16, 46}
\definecolor{valdark}       {RGB}{ 40, 40, 40}
\definecolor{valorange}     {RGB}{230,100,  0}
\definecolor{valpurple}     {RGB}{120, 70,160}
\definecolor{vallightblue}  {RGB}{220,235,250}
\definecolor{vallightgreen} {RGB}{220,245,230}
\definecolor{vallightred}   {RGB}{255,228,228}
\definecolor{vallightorange}{RGB}{255,238,210}
\definecolor{valgray}       {RGB}{245,245,245}

% ═══════════════════════════════════════════════════════════════════════
%  AMBIENTES TCOLORBOX INSTITUCIONALES
% ═══════════════════════════════════════════════════════════════════════
\newtcolorbox{concepto}[1]{
	enhanced,breakable, title={\large\bfseries\ding{111} Aprende el Concepto: #1},
	colback=vallightblue!40, colframe=valblue, coltitle=white,
	attach boxed title to top left={yshift=-2mm,xshift=4mm},
	boxed title style={colback=valblue,colframe=valblue,rounded corners},
	arc=4mm, boxrule=1.5pt, top=6mm, left=3mm, right=3mm,
}
\newtcolorbox{truco}[1]{
	enhanced,breakable, title={\large\bfseries\ding{51} El Truco: #1},
	colback=vallightorange!40, colframe=valorange, coltitle=white,
	attach boxed title to top left={yshift=-2mm,xshift=4mm},
	boxed title style={colback=valorange,colframe=valorange,rounded corners},
	arc=4mm, boxrule=1.5pt, top=6mm, left=3mm, right=3mm,
}
\newtcolorbox{atencion}[1]{
	enhanced,breakable, title={\large\bfseries\ding{55} #1},
	colback=vallightred!40, colframe=valred, coltitle=white,
	attach boxed title to top left={yshift=-2mm,xshift=4mm},
	boxed title style={colback=valred,colframe=valred,rounded corners},
	arc=4mm, boxrule=1.5pt, top=6mm, left=3mm, right=3mm,
}
\newtcolorbox{practica}[1]{
	enhanced,breakable, title={\large\bfseries\ding{46} ¡Tu Turno! (#1)},
	colback=white, colframe=valdark, coltitle=white,
	attach boxed title to top left={yshift=-2mm,xshift=4mm},
	boxed title style={colback=valdark,colframe=valdark,rounded corners},
	arc=4mm, boxrule=1.5pt, top=6mm, left=3mm, right=3mm,
}

% ── Casilla de respuesta ajustada ───────────────────────
\newcommand{\boxfill}[1][2.2cm]{%
	\tikz[baseline=-0.5ex]%
	\draw[valdark!60, fill=valgray, line width=0.8pt, rounded corners=1pt] (0,-0.3) rectangle (#1,0.3cm);%
}

% ── Encabezado de Misión ──────────────────────────────────────────────
\newcommand{\mision}[2]{%
	\needspace{4cm}%
	\vspace{4pt}%
	\begin{tcolorbox}[
		enhanced, colback=#1, colframe=#1!60!black, arc=3mm, boxrule=1.5pt,
		left=6pt, right=6pt, top=4pt, bottom=4pt,
		]
		{\Large\bfseries\color{white} #2}
	\end{tcolorbox}%
	\vspace{2pt}%
}

\pagestyle{plain}
\setlength{\parindent}{0pt}
\setlength{\parskip}{2pt}

% ============================================================
\begin{document}
	
	% ─────────────────────────────────────────────────────────────
	% PORTADA
	% ─────────────────────────────────────────────────────────────
	\begin{center}
		{\Large\bfseries\color{valblue} LICEO V.A.L.\ (VIDA, AMOR, LUZ)}\\[4pt]
		{\large\bfseries MÓDULO DE AUTOAPRENDIZAJE\ ---\ MATEMÁTICAS 6\textdegree}\\[8pt]
		\begin{tcolorbox}[enhanced,arc=6mm,boxrule=2pt,
			colback=valblue,colframe=valdark,halign=center]
			{\Huge\bfseries\color{white} UNIDAD 17}\\[4pt]
			{\Large\bfseries\color{yellow!90!white} Estadística. Gráfica de Barras.}\\[2pt]
			{\large\bfseries\color{white} Lectura de tablas, gráficas y pictogramas}
		\end{tcolorbox}
		\vspace{4pt}
		\begin{tcolorbox}[enhanced,arc=4mm,boxrule=1.5pt,
			colback=white,colframe=valblue,width=0.88\linewidth,halign=center]
			{\bfseries Nombre:}\\[6pt]
			\rule{9cm}{0.5pt}\\[6pt]
			{\bfseries Código:}\ \textbf{06M-17 FL}
			\hspace{2cm}
			{\bfseries Año:}\ \underline{2026}
		\end{tcolorbox}
		\vspace{2pt}
		{\itshape ``Hacia un Aprendizaje Significativo, Autónomo y en Libertad''}
	\end{center}
	\vspace{6pt}
	
	% =========================================================
	%  MISIÓN 1 — TIPOS DE GRÁFICAS Y LECTURA DE TABLAS
	% =========================================================
	\mision{valblue}{Misión 1: Tipos de Gráficas y Lectura de Tablas}
	
	\begin{concepto}{La Estadística y sus Representaciones}
		La \textbf{estadística} recoge, organiza e interpreta datos para sacar conclusiones. Un mismo conjunto de datos puede representarse de \textbf{varias formas} según lo que se quiera mostrar.
		
		\vspace{4pt}
		\renewcommand{\arraystretch}{1.5}
		\begin{center}
			\small
			\begin{tabularx}{\linewidth}{|>{\bfseries\color{valblue}}m{3.5cm}|X|}
				\hline
				\rowcolor{valblue!20}
				\textcolor{valdark}{Tipo de gráfica} & \textbf{¿Para qué se usa?} \\
				\hline
				Gráfica de barras & Compara \textbf{cantidades} entre distintas categorías usando rectángulos. \\
				\hline
				Gráfica circular (torta) & Muestra cómo se \textbf{reparte un total} en partes proporcionales (porcentajes). \\
				\hline
				Pictograma & Usa \textbf{símbolos o íconos} repetidos. Cada símbolo es una cantidad fija. \\
				\hline
				Gráfica de líneas & Conecta puntos para mostrar la \textbf{evolución} de un dato en el \textbf{tiempo}. \\
				\hline
				Histograma & Agrupa datos numéricos \textbf{continuos} en intervalos (barras unidas sin espacio).\\
				\hline
			\end{tabularx}
		\end{center}
	\end{concepto}
	
	\begin{atencion}{¡Cuidado! Barras vs. Histograma}
		En la gráfica de barras hay \textbf{espacio} entre cada barra porque representa \textbf{categorías distintas} (países, colores). En el histograma las barras están \textbf{pegadas} porque representan \textbf{intervalos consecutivos} de un número (edades de 0-10, 10-20...).
	\end{atencion}
	
	\begin{truco}{¿Cómo leer una Tabla de Datos paso a paso?}
		Para encontrar información precisa en una tabla sin perderte, imagina que juegas a hundir la flota o al ajedrez, cruzando coordenadas.
		\begin{itemize}[itemsep=2pt, leftmargin=*]
			\item \textbf{Paso 1 (Encabezados):} Revisa qué significa cada fila (horizontal) y cada columna (vertical).
			\item \textbf{Paso 2 (La Intersección):} El dato que buscas está exactamente donde se \textbf{cruza} la fila con la columna.
		\end{itemize}
		\textbf{Ejemplo:} Si te preguntan \textit{"¿Cuántos Goles hizo Mateo en el Partido 2?"}, buscas la fila \textbf{Mateo}, la columna \textbf{Partido 2}, e intersectas ambas líneas. ¡El resultado es \textbf{2}!
		
		\begin{center}
			\small
			\renewcommand{\arraystretch}{1.2}
			\begin{tabular}{|l|c|c|c|}
				\hline
				\rowcolor{valorange!20} \textbf{Jugador} & \textbf{Partido 1} & \textbf{\textcolor{valred}{Partido 2}} & \textbf{Partido 3} \\ \hline
				Lucas & 1 & 0 & 3 \\ \hline
				\rowcolor{vallightred!40} \textbf{\textcolor{valred}{Mateo}} & 0 & \textbf{\textcolor{valred}{2}} $\leftarrow$ & 1 \\ \hline
				Simón & 2 & 1 & 0 \\ \hline
			\end{tabular}
		\end{center}
	\end{truco}
	
	\begin{practica}{Tipos de Gráficas y Lectura de Tablas}
		\textbf{A) Identifica el tipo de gráfica según la descripción:}
		\vspace{-4pt}
		\begin{multicols}{2}
			\begin{enumerate}[label=\textbf{\arabic*.}, itemsep=0.2cm, leftmargin=*]
				\small
				\item Usa símbolos/íconos repetidos: \boxfill[2.5cm]
				\item Divide un círculo en porcentajes: \boxfill[2.5cm]
				\item Barras separadas para categorías: \boxfill[2.5cm]
				\item Puntos y líneas para el tiempo: \boxfill[2.5cm]
				\item Barras unidas para intervalos: \boxfill[2.5cm]
			\end{enumerate}
		\end{multicols}
		
		\vspace{6pt}
		\hrule
		\vspace{6pt}
		
		\textbf{B) Análisis de Tablas:} Identifica qué representa cada fila y columna.
		
		\begin{minipage}[t]{0.45\linewidth}
			\centering
			\vspace{0pt}
			\textbf{Vuelos cancelados por aerolínea}\\[4pt]
			\renewcommand{\arraystretch}{1.5}
			\small
			\begin{tabular}{|l|c|c|c|c|}
				\hline
				\rowcolor{valblue!20}
				\textbf{Aerolínea} & \textbf{Ene} & \textbf{Feb} & \textbf{Mar} & \textbf{Abr} \\
				\hline
				AeroSur     & 12 & 8  & 15 & 10 \\ \hline
				Vuela Alto  & 5  & 9  & 6  & 14 \\ \hline
				CieloAzul   & 18 & 3  & 7  & 9  \\ \hline
				TransAndes  & 9  & 12 & 4  & 12 \\ \hline
			\end{tabular}
		\end{minipage}%
		\hfill
		\begin{minipage}[t]{0.52\linewidth}
			\vspace{0pt}
			\begin{enumerate}[start=6, itemsep=0.2cm, leftmargin=0.5cm, label=\textbf{\arabic*.}]
				\small
				\item Vuelos cancelados por \textbf{AeroSur} en marzo:\quad\boxfill[1.2cm]
				\item Vuelos cancelados por \textbf{CieloAzul} en enero:\quad\boxfill[1.2cm]
				\item Aerolínea con \textbf{más} cancelaciones en abril:\quad\boxfill[2cm]
				\item Aerolínea con \textbf{menos} cancelaciones en febrero:\quad\boxfill[2cm]
				\item Mes con \textbf{más} cancelaciones para \textbf{Vuela Alto}:\quad\boxfill[2cm]
			\end{enumerate}
		\end{minipage}
		
		\vspace{8pt}
		\begin{enumerate}[start=11, itemsep=0.3cm, leftmargin=0.5cm, label=\textbf{\arabic*.}]
			\small
			\item ¿Cuántos vuelos canceló \textbf{TransAndes en total} en los 4 meses?  \boxfill[1.5cm]
			\item ¿Cuántos vuelos \textbf{más} canceló CieloAzul en enero que en febrero?  \boxfill[3.5cm]
			\item ¿Qué aerolínea canceló el \textbf{mismo número} de vuelos en dos meses distintos? ¿En cuáles? \boxfill[5cm]
			\item Sumando las 4 aerolíneas, ¿cuántos vuelos se cancelaron en \textbf{enero}?  \boxfill[1.5cm]
			\item ¿Cuál fue el \textbf{mes con mayor total} de cancelaciones (sumando todas)?  \boxfill[3.5cm]
		\end{enumerate}
	\end{practica}
	
	%\newpage
	% =========================================================
	%  MISIÓN 2 — ANÁLISIS PROFUNDO DE GRÁFICAS DE BARRAS
	% =========================================================
	\mision{valorange}{Misión 2: Análisis de Gráficas de Barras (Positivos y Negativos)}
	
	\begin{concepto}{Lectura con Línea de Cero}
		Una gráfica de barras puede mostrar \textbf{números enteros}. Las barras que suben desde la línea del cero representan \textbf{ganancias} (+) y las que bajan representan \textbf{pérdidas} (-). El cero es tu nivel de suelo o punto de referencia.
	\end{concepto}
	
	\begin{truco}{¿Cómo deducir valores que no están escritos?}
		En gráficas avanzadas los valores \textbf{no se escriben sobre la barra}. Debes deducirlos así:
		\begin{enumerate}[itemsep=2pt]
			\small
			\item Mira el \textbf{tope} de la barra (su punto más alto o más bajo).
			\item Traza una \textbf{línea imaginaria horizontal} recta hacia el eje numérico (Eje Y).
			\item Lee el número. Si cae entre dos líneas, calcula el valor intermedio basándote en la escala.
		\end{enumerate}
		\vspace{4pt}
		\centering
		\begin{tikzpicture}[x=1.2cm,y=0.45cm]
			\draw[gray!30,line width=0.5pt] (0,0) grid[xstep=4, ystep=2] (3,8);
			\foreach \y in {0,2,4,6,8} { \node[left,font=\bfseries\scriptsize,color=valdark] at (0,\y) {\y}; }
			\draw[valdark,line width=1.2pt] (0,0) -- (0,9);
			\draw[valdark,line width=1.2pt] (0,0) -- (3.5,0);
			
			% Barra ejemplo
			\path[fill=valblue!75, draw=valblue, line width=1pt] (1.5,0) rectangle (2.5,6);
			\node[below, font=\scriptsize] at (2.0,0) {Categoría};
			
			% Línea guía visual
			\draw[valorange, dashed, line width=1.5pt, ->, >=stealth] (1.5, 6) -- (0.2, 6);
			\fill[valorange] (0,6) circle (2.5pt);
			\node[font=\bfseries\scriptsize, color=valorange, fill=white, inner sep=2pt] at (1.2, 6.7) {¡Sigue esta línea!};
		\end{tikzpicture}
	\end{truco}
	
	\begin{practica}{Ganancias y Pérdidas - Constructora ANDES}
		
		\begin{minipage}[t]{0.48\linewidth}
			\centering
			\vspace{0pt}
			\begin{tikzpicture}[x=0.9cm,y=0.12cm]
				% Cuadrícula de fondo cada 5 unidades para fácil lectura
				\foreach \y in {-30,-25,-20,-15,-10,-5,0,5,10,15,20,25,30,35,40,45,50}{
					\draw[gray!20, line width=0.5pt] (0,\y) -- (7.5,\y);
				}
				% Líneas principales cada 10 unidades
				\foreach \y in {-30,-20,-10,0,10,20,30,40,50}{
					\draw[gray!60, dashed, line width=0.6pt] (0,\y) -- (7.5,\y);
					\node[left, font=\bfseries\scriptsize, color=valdark] at (0,\y) {\y};
				}
				% Eje Y y Eje X (Cero)
				\draw[valdark, line width=1.5pt] (0,-35) -- (0,55);
				\draw[valdark, line width=1.8pt] (0,0) -- (7.8,0);
				
				\node[rotate=90, font=\bfseries\small, valdark] at (-1.2,10) {Millones de \$};
				\node[font=\bfseries\small, valorange, align=center, text width=6cm] at (3.7,63) {Ganancias y Pérdidas (2007 - 2012)};
				
				% Barras (año: valor en millones) - SIN VALORES EXPLÍCITOS
				\foreach \x/\y/\col/\name in {
					1.0/20/valgreen/2007,
					2.2/35/valgreen/2008,
					3.4/-15/valred/2009,
					4.6/-25/valred/2010,
					5.8/10/valgreen/2011,
					7.0/45/valgreen/2012%
				} {
					\path[fill=\col!75, draw=\col, line width=1pt] (\x-0.4,0) rectangle (\x+0.4,\y);
					% Condicional para poner el año arriba o abajo dependiendo si es pérdida o ganancia
					\ifnum\y<0
					\node[above=2pt, font=\bfseries\scriptsize, color=valdark] at (\x,0) {\name};
					\else
					\node[below=2pt, font=\bfseries\scriptsize, color=valdark] at (\x,0) {\name};
					\fi
				}
			\end{tikzpicture}
		\end{minipage}%
		\hfill
		\begin{minipage}[t]{0.50\linewidth}
			\vspace{0pt}
			\begin{enumerate}[itemsep=0.35cm, leftmargin=0.5cm, label=\textbf{\arabic*.}]
				\small
				\item Ganancia en \textbf{2007}: \hfill\boxfill[2cm]
				\item Pérdida en \textbf{2009}: \hfill\boxfill[2cm]
				\item Año con \textbf{mayor ganancia}: \hfill\boxfill[2cm]
				\item Año con \textbf{mayor pérdida}: \hfill\boxfill[2cm]
				\item ¿Cuántos años tuvieron \textbf{pérdidas}? \hfill\boxfill[1cm]
				\item ¿Cuántos años tuvieron \textbf{ganancias}? \hfill\boxfill[1cm]
				\item Suma las ganancias (2007, 2008, 2011, 2012): \hfill\boxfill[2cm]
				\item Suma las pérdidas (2009, 2010): \hfill\boxfill[2cm]
			\end{enumerate}
		\end{minipage}
		
		\vspace{10pt}
		\hrule
		\vspace{10pt}
		
		\begin{enumerate}[start=9, itemsep=0.4cm, leftmargin=0.5cm, label=\textbf{\arabic*.}]
			\small
			\item Calcula la \textbf{diferencia neta} entre 2008 y 2009 (resta el valor de 2009 al de 2008, cuidando el signo). \hfill Resultado: \boxfill[2cm]
			\item Calcula el \textbf{balance total} sumando todos los años (ganancias positivas y pérdidas negativas). \hfill Balance: \boxfill[2.5cm]
			\item Compara 2010 y 2011: ¿cuántos millones \textbf{mejoró} el resultado de la empresa de un año al otro? \hfill \boxfill[2cm]
			\item Si en 2013 la empresa gana 10 millones \textbf{más} que en 2012, ¿cuánto ganaría en 2013? \hfill \boxfill[2cm]
			\item Ordena de \textbf{menor a mayor} los 6 valores de la gráfica: \hfill \boxfill[8cm]
		\end{enumerate}
	\end{practica}
	
	\newpage
	% =========================================================
	%  MISIÓN 3 — GRÁFICAS DE LÍNEAS
	% =========================================================
	\mision{valgreen}{Misión 3: Gráficas de Líneas (Tendencias Temporales)}
	
	\begin{concepto}{Lectura de una Gráfica de Líneas}
		El eje \textbf{horizontal (X)} representa el \textbf{tiempo}, y el eje \textbf{vertical (Y)} la cantidad que cambia. La línea puede:
		\textbf{Subir} (aumento), \textbf{Bajar} (disminución), o mantenerse \textbf{Horizontal} (constante). Para leer valores exactos, ubica el punto en la línea y usa el \textbf{Truco de la línea imaginaria} hacia la izquierda (eje Y).
	\end{concepto}
	
	\begin{practica}{Altitud de un avión durante un vuelo}
		
		\begin{minipage}[t]{0.52\linewidth}
			\centering
			\vspace{0pt}
			\begin{tikzpicture}[x=1.1cm,y=0.18cm]
				% Cuadrícula cada 2 unidades para exactitud
				\foreach \y in {0,2,4,6,8,10,12,14,16,18,20,22,24,26,28,30,32,34,36}{
					\draw[gray!20, line width=0.4pt] (0,\y) -- (7.5,\y);
				}
			
				
				% --- Cuadrícula y números del Eje Y de 2 en 2 ---
				\foreach \y in {2, 4, 6, 8, 10, 12, 14, 16, 18, 20, 22, 24, 26, 28, 30, 32, 34, 36} {
					% 1. Línea horizontal de la cuadrícula (gris clarito)
					\draw[gray!22, line width=0.55pt] (0,\y) -- (8.5,\y);
					
					% 2. Pequeña marca sobre el eje Y principal
					\draw[valdark] (-0.12,\y) -- (0,\y);
					
					% 3. El número al lado izquierdo
					\node[left, font=\bfseries\tiny] at (0,\y) {\y};
				}
				% Etiqueta manual para el cero en el origen
				\node[left, font=\bfseries\tiny] at (0,0) {0};
				
				% Ejes
				\draw[valdark, line width=1.5pt] (0,0) -- (0,38);
				\draw[valdark, line width=1.5pt] (0,0) -- (7.8,0);
				
				\node[rotate=90, font=\bfseries\small, valdark] at (-1.0,18) {Altitud (miles de pies)};
				\node[font=\bfseries\small, valdark] at (3.5,-3.8) {Horas de vuelo};
				\node[font=\bfseries\small, valgreen, align=center] at (3.5,42) {Altitud del Vuelo AV-2026};
				
				% Eje X marcas
				\foreach \x in {0,1,2,3,4,5,6,7}{
					\draw[valdark, line width=1pt] (\x,0) -- (\x,-1);
					\node[below, font=\bfseries\scriptsize] at (\x,-1) {\x};
				}
				
				% Línea y puntos (Sin etiquetas)
				\draw[valblue, line width=2pt] (0,0) -- (1,15) -- (2,30) -- (3,32) -- (4,32) -- (5,32) -- (6,18) -- (7,0);
				
				\foreach \x/\y in {0/0, 1/15, 2/30, 3/32, 4/32, 5/32, 6/18, 7/0}{
					\filldraw[valred, draw=white, line width=1pt] (\x,\y) circle (2.5pt);
				}
			\end{tikzpicture}
		\end{minipage}%
		\hfill
		\begin{minipage}[t]{0.45\linewidth}
			\vspace{0pt}
			\begin{enumerate}[itemsep=0.4cm, leftmargin=0.5cm, label=\textbf{\arabic*.}]
				\small
				\item Altitud en la \textbf{hora 1}: \boxfill[2cm]
				\item Altitud en la \textbf{hora 2}: \boxfill[2cm]
				\item Altitud en la \textbf{hora 4}: \boxfill[2cm]
				\item Altitud en la \textbf{hora 6}: \boxfill[2cm]
				\item \textbf{Máxima altitud} alcanzada: \boxfill[2cm]
				\item ¿Cuántas \textbf{horas duró} el vuelo? \boxfill[1.5cm]
				\item Intervalo de horas donde la altitud fue \textbf{constante}: \boxfill[2cm]
				\item Intervalo de horas donde el avión estuvo \textbf{descendiendo}: \boxfill[2cm]
			\end{enumerate}
		\end{minipage}
		
		\vspace{10pt}
		\hrule
		\vspace{10pt}
		
		\begin{enumerate}[start=9, itemsep=0.4cm, leftmargin=0.5cm, label=\textbf{\arabic*.}]
			\small
			\item Calcula cuánto \textbf{ascendió} el avión entre la hora 0 y la hora 1.  \boxfill[2.5cm]
			\item Calcula cuánto \textbf{descendió} el avión entre la hora 5 y la hora 6. \boxfill[2.5cm]
			\item Compara el ascenso entre las horas 0-1, con el ascenso entre las horas 1-2. ¿Cuál fue mayor?  \boxfill[3cm]
			\item ¿Cuál fue la \textbf{diferencia} entre la altitud máxima y la altitud en la hora 6?  \boxfill[2.5cm]\\
			\newpage
			\item En tus palabras, describe \textbf{la historia completa} del vuelo usando ``sube'', ``se mantiene'' y ``baja''.\\
			\rule{\linewidth}{0.5pt}\\
			\rule{\linewidth}{0.5pt}
		\end{enumerate}
	\end{practica}
	
	% =========================================================
	%  MISIÓN 4 — GRÁFICAS CIRCULARES Y PICTOGRAMAS
	% =========================================================
	\mision{valred}{Misión 4: Gráficas Circulares y Pictogramas}
	
	\begin{concepto}{Proporciones y Símbolos}
		En una \textbf{gráfica circular}, el círculo completo es el 100\%. Para hallar los elementos de una porción, multiplica el \textbf{porcentaje (en decimal)} por el \textbf{total} (Ej: 25\% es 0.25). 
		En un \textbf{pictograma}, cada símbolo vale una \textbf{cantidad fija}. Multiplica la cantidad de símbolos por ese valor. Medio símbolo es la mitad del valor.
	\end{concepto}
	
	\begin{practica}{A) Gráfica Circular (Total: 240 estudiantes)}
		
		\begin{minipage}[c]{0.45\linewidth}
			\centering
			\vspace{0pt}
			\begin{tikzpicture}[scale=1.1]
				% Deporte: 144°, 40% (90 a -54)
				\draw[fill=valblue!75, draw=white, line width=1pt] (0,0) -- (90:2.5) arc (90:-54:2.5) -- cycle;
				% Lectura: 90°, 25% (-54 a -144)
				\draw[fill=valgreen!75, draw=white, line width=1pt] (0,0) -- (-54:2.5) arc (-54:-144:2.5) -- cycle;
				% Música: 72°, 20% (-144 a -216)
				\draw[fill=valorange!75, draw=white, line width=1pt] (0,0) -- (-144:2.5) arc (-144:-216:2.5) -- cycle;
				% Arte: 54°, 15% (-216 a -270)
				\draw[fill=valred!75, draw=white, line width=1pt] (0,0) -- (-216:2.5) arc (-216:-270:2.5) -- cycle;
				
				\node[font=\bfseries\small, white] at (18:1.5) {40\%};
				\node[font=\bfseries\small, white] at (-99:1.6) {25\%};
				\node[font=\bfseries\small, white] at (180:1.6) {20\%};
				\node[font=\bfseries\small, white] at (117:1.6) {15\%};
				
				% Leyenda
				\begin{scope}[shift={(3.2, 1.2)}]
					\fill[valblue!75] (0,0) rectangle (0.4,0.4); \node[right, font=\small] at (0.5,0.2) {Deporte};
					\fill[valgreen!75] (0,-0.7) rectangle (0.4,-0.3); \node[right, font=\small] at (0.5,-0.5) {Lectura};
					\fill[valorange!75] (0,-1.4) rectangle (0.4,-1.0); \node[right, font=\small] at (0.5,-1.2) {Música};
					\fill[valred!75] (0,-2.1) rectangle (0.4,-1.7); \node[right, font=\small] at (0.5,-1.9) {Arte};
				\end{scope}
			\end{tikzpicture}
		\end{minipage}%
		\hfill
		\begin{minipage}[c]{0.50\linewidth}
			\vspace{0pt}
			\begin{enumerate}[itemsep=0.4cm, leftmargin=0.5cm, label=\textbf{\arabic*.}]
				\small
				\item Estudiantes en \textbf{Deporte} (40\% de 240): \hfill\boxfill[2cm]
				\item Estudiantes en \textbf{Lectura} (25\% de 240): \hfill\boxfill[2cm]
				\item Estudiantes en \textbf{Música} (20\% de 240): \hfill\boxfill[2cm]
				\item Estudiantes en \textbf{Arte} (15\% de 240): \hfill\boxfill[2cm]
				\item ¿Cuántos estudiantes \textbf{más} prefieren Deporte que Arte? \hfill\boxfill[2cm]
				\item ¿Qué \textbf{fracción} simplificada representa Música (20\%)? \hfill\boxfill[2cm]
			\end{enumerate}
		\end{minipage}
	\end{practica}
	
	\begin{practica}{B) Pictograma: Mascotas de grado 6º}
		
		\begin{minipage}[c]{0.50\linewidth}
			\centering
			\vspace{0pt}
			\textit{Cada símbolo\; \tikz[baseline=-2pt]{\fill[valblue!75, rounded corners=1pt] (0,0) rectangle (0.35,0.35);}\; vale \textbf{4 mascotas}.}\\
			\vspace{6pt}
			\renewcommand{\arraystretch}{1.8}
			\begin{tabular}{|l|l|}
				\hline
				\rowcolor{valblue!15}
				\textbf{Mascota} & \textbf{Símbolos} \\ \hline
				Perros &
				\begin{tikzpicture}[baseline=-2pt]
					\foreach \i in {0,1,2,3,4}{ \fill[valblue!75, rounded corners=1pt] (\i*0.45,0) rectangle (\i*0.45+0.35,0.35); }
				\end{tikzpicture} \\ \hline
				Gatos &
				\begin{tikzpicture}[baseline=-2pt]
					\foreach \i in {0,1,2}{ \fill[valblue!75, rounded corners=1pt] (\i*0.45,0) rectangle (\i*0.45+0.35,0.35); }
					\fill[valblue!75, rounded corners=1pt] (1.35,0) rectangle (1.525,0.35); % Mitad
				\end{tikzpicture} \\ \hline
				Peces &
				\begin{tikzpicture}[baseline=-2pt]
					\foreach \i in {0,1}{ \fill[valblue!75, rounded corners=1pt] (\i*0.45,0) rectangle (\i*0.45+0.35,0.35); }
				\end{tikzpicture} \\ \hline
				Aves &
				\begin{tikzpicture}[baseline=-2pt]
					\fill[valblue!75, rounded corners=1pt] (0,0) rectangle (0.35,0.35);
				\end{tikzpicture} \\ \hline
			\end{tabular}
		\end{minipage}%
		\hfill
		\begin{minipage}[c]{0.45\linewidth}
			\vspace{0pt}
			\begin{enumerate}[start=7, itemsep=0.4cm, leftmargin=0.5cm, label=\textbf{\arabic*.}]
				\small
				\item Mascotas tipo \textbf{Perro}: \hfill\boxfill[2cm]
				\item Mascotas tipo \textbf{Gato}: \hfill\boxfill[2cm]
				\item Mascotas tipo \textbf{Pez}: \hfill\boxfill[2cm]
				\item Mascotas tipo \textbf{Ave}: \hfill\boxfill[2cm]
				\item Mascotas \textbf{en total} (sumando las 4): \hfill\boxfill[2cm]
				\item ¿Cuántas mascotas \textbf{más} hay de tipo Perro que de tipo Ave? \hfill\boxfill[2cm]
			\end{enumerate}
		\end{minipage}
	\end{practica}
	
\end{document}